\chapter{Referencia rápida}

\begin{lstlisting}[style=bash]
# Preparar la terminal
source /opt/ros/jazzy/setup.bash && source ~/ros2_ws/install/setup.bash

# Búsqueda autónoma con el STA / con ArduPilot
ros2 launch cuadrirotor_sim simulacion.launch.py vehiculo:=tarot650 mundo:=bosque_grande controlador:=sta mision:=busqueda
ros2 launch cuadrirotor_sim simulacion.launch.py vehiculo:=tarot650 mundo:=bosque_grande controlador:=ardupilot mision:=busqueda

# Aviso del incendio y estado
ros2 topic echo /cuadrirotor/incendio/alerta
ros2 topic echo /cuadrirotor/busqueda/estado
cat ~/incendios_detectados.csv

# Pruebas
cd ~/ros2_ws/src/cuadrirotor_sim && python3 -m pytest -q test/
\end{lstlisting}

\chapter{Configuración del mundo grande}

\cod{config/mundo_bosque_grande.yaml}:
\begin{lstlisting}
# Búsqueda en el mundo 'bosque_grande' (valle de 420 x 420 m con lomas y cerros).
# El launch agrega 'terreno' (el DEM); las alturas son sobre el suelo.
busqueda_incendio:
  ros__parameters:
    zona: [-210.0, 210.0, -210.0, 210.0]   # xmin, xmax, ymin, ymax del bosque [m]
    patron: serpentina        # avanza hacia el frente del dron (franjas: boustrophedon)
    separacion: 70.0          # entre barridos [m]
    altura_busqueda: 30.0     # sobre el suelo; los árboles miden hasta ~8 m
    anticipacion: 50.0        # mira el relieve esta distancia hacia adelante [m]
    velocidad: 5.0
    velocidad_vertical: 3.0
    aceleracion: 1.5
    alcance_vision: 150.0     # para el mapa de cobertura y el corte del rayo [m]
    # Órbita por encima de las copas (el planificador no sabe dónde hay árboles).
    altura_orbita: 18.0
    altura_minima_copas: 14.0
    radio: 14.0
    omega: 0.12
    # Evasión: no avanza a menos de 10 m sobre el relieve que tiene adelante
    # (primero sube); el LiDAR rodea lo que esté a menos de 3 m del camino.
    altura_seguridad: 10.0
    radio_evasion: 3.0
    alcance_evasion: 20.0
\end{lstlisting}

\chapter{YAML del Tarot generado}

\cod{config/vehiculo_tarot650.yaml} (lo escribe \cod{generar_modelo.py}):
\begin{lstlisting}
/**:
  ros__parameters:
    masa: 2.2
    inercia: [0.03412, 0.03609, 0.06596]
    rotores_xy: [0.2198, -0.212, 0.2198, 0.212, -0.2198, 0.212, -0.2198, -0.212]
    rotores_sentido: [1, -1, 1, -1]
    kf: 4.24e-05
    km: 0.016
    w_max: 680.0
    kr_att: [3.4, 3.6, 1.65]
    kw_att: [0.55, 0.58, 0.59]
    z_inicial: 0.235
    camara_xyz: [0.12, 0.0, -0.06]
    camara_pitch: 0.5236
    camara_hfov: 1.3
    camara_ancho: 640
    camara_alto: 480

controlador_sta:
  ros__parameters:
    mu: [0.3, 0.3, 0.03]
    k1_max: 3.0
\end{lstlisting}

\chapter{Créditos de los recursos}

\begin{itemize}
  \item Tarot Ironman 650: modelo CAD de GrabCAD,
        \url{https://grabcad.com/library/tarot-ironman-650-2}.
  \item Pino y roble: OpenRobotics en Gazebo Fuel (CC0).
  \item Venado (crownjoshua), zorro y lobo (Quaternius): OpenGameArt, CC0.
        \url{https://opengameart.org/content/cc0-3d-animals-creatures},
        \url{https://opengameart.org/content/deer-low-poly-rigged},
        \url{https://opengameart.org/content/5-low-poly-animals},
        \url{https://opengameart.org/content/animated-animales-low-poly}.
  \item YOLO26: Ultralytics. ArduPilot y \cod{ardupilot_gazebo}: proyecto ArduPilot (GPLv3).
  \item Imagen satelital del mapa de MAVProxy: la que descarga MAVProxy para la ubicación
        supuesta.
\end{itemize}
